\section{Feasible transport under endogenous constraints}\label{app:feasible46}
We give the details underlying the feasible-witness transport theorem in the main article. All norms below are the specified $\ell^1$ norms. The evaluation functional is the same at the compared state--action pairs. Its monotonicity and cash invariance imply that changing its bounded random input by at most $b$ changes its value by at most $b$. Thus the primitive state and action inequalities on the feasible graph imply
\[
 |Q_t(f_{t+1};x,a)-Q_t(f_{t+1};z,b)|
 \le A_t^Q\|x-z\|_1+D_t^Q\|a-b\|_1
\]
whenever $a\in A_t(x)$ and $b\in A_t(z)$. The proof never evaluates $Q_t$ outside that graph.

\subsection{The Bellman modulus and noisy feasible labels}
Fix $a\in A_t(x)$. Compactness and the Hausdorff condition give $b\in A_t(z)$ with $\|a-b\|_1\le K_t^A\|x-z\|_1$. It follows that
\[
 \mathcal T_tf_{t+1}(z)
 \le Q_t(f_{t+1};x,a)+L_t\|x-z\|_1.
\]
Taking the infimum over $a$ and then reversing $x,z$ establishes the claimed Bellman modulus, without assuming that the minimizer is known. At node $i$, the minimum of exact objective values over its feasible action net is at least the continuous feasible infimum and at most that infimum plus $D_t^Qk_t$. Taking the minimum of queries with absolute error at most $e_t$ adds at most $e_t$ in each direction.

For every cone, the lower label inequality and the Bellman modulus give
\[
 y_{t,i}+L_t\|x-x_{t,i}\|_1\ge\mathcal T_tf_{t+1}(x)-e_t.
\]
For a nearest state node, the upper label inequality gives an upper bound $\mathcal T_tf_{t+1}(x)+D_t^Qk_t+e_t+2L_th_t$. The finite minimum has the same bounds and is $L_t$-Lipschitz. The next backward step therefore has precisely the modulus required by the construction, independently of the numerical label errors.

\subsection{Repair, selection and the dynamic account}
The stored action $a_{t,i}$ attains an original numerical query; hence its exact node objective is at most $y_{t,i}+e_t$. For the minimizing cone, the feasible-pair inequality and repair displacement give
\[
 Q_t(f_{t+1};x,\Gamma_{t,i}(x))-f_t(x)
 \le e_t+D_t^Q\zeta_t.
\]
A fixed ordering of finitely many continuous cone values gives a Borel index selector. Composing the finitely many Borel repairs with that selector gives a feasible Borel policy. The one-sided sandwich now supplies the discounted loss bound, with the original terminal interval unchanged. The full-policy comparison is exactly the Bellman comparison assumed in the controlled economy; it is not a comparison restricted to the node action arrays.

The four nonterminal resource allocations make the bracket at most $s/4+s/4+s/4+s/4=s$. The terminal bracket is at most $s$. Effective finite covers, feasible action nets, query enclosures and repair procedures consequently give finite termination. Where only an approximate repair is computable, its certified displacement is charged through $D_t^Q\zeta_t$ and its actual action must still be feasible. A numerical tolerance for an infeasible constraint violation is not included in this theorem.

\subsection{A transparent exact check}
Let $X=[0,1]^2$, $b(x)=(x_1+x_2)/2$, $A(x)=[0,b(x)]$, and consider the one-date objective
\[
 Q(x,a)=(a-3/4)^2+x_1/8+x_2/16.
\]
Here $A^Q=1/8$, $D^Q=3/2$, $K^A=1/2$, and $L=7/8$ are valid moduli. The exact optimum is attained at $\min\{3/4,b(x)\}$. For a feasible anchor action, clipping at $b(x)$ changes the action by at most $|b(x)-b(x_i)|\le\|x-x_i\|_1/2$. The transport theorem therefore applies directly, and a reference implementation can compare the exact repaired-policy loss with the analytic bound using rational arithmetic.

The additional tests independently form feasible node nets and noisy labels, compute the cone minimum by direct enumeration, and check the Bellman and selected-policy enclosures on a fixed rational set. They also check the Hausdorff modulus, exact feasibility, approximate-repair accounting, nonminimizing-index allowances and the zero-constraint-variation reduction. The analytic proof above, rather than a finite grid test, establishes the all-state claim. The tests are not training data, timing observations, or evidence of a high-dimensional speed advantage. No R46 catalogue record is altered by these checks.
